% RetrievalBench paper draft
% Compile: pdflatex main.tex && bibtex main && pdflatex main.tex && pdflatex main.tex
\documentclass[10pt]{article}
\usepackage[margin=1in]{geometry}
\usepackage{booktabs}
\usepackage{amsmath,amssymb}
\usepackage{graphicx}
\usepackage{hyperref}
\usepackage{multirow}
\usepackage{array}
\usepackage{xcolor}
\usepackage{caption}

\title{RetrievalBench: Production-Realistic Evaluation \\
Shows BEIR Rankings Do Not Predict RAG Retrieval Performance}
\author{Anonymous}
\date{}

\begin{document}
\maketitle

\begin{abstract}
BEIR established zero-shot cross-domain retrieval as the evaluation paradigm of the neural era, but production Retrieval-Augmented Generation (RAG) introduces failure modes that BEIR was never designed to measure: chunking artifacts, context-window saturation, latency SLAs, and per-query cost budgets. We introduce \textbf{RetrievalBench}, an evaluation framework that treats chunking strategy, reranking, latency, and cost as first-class variables alongside the retriever itself. Across seven BEIR domains spanning scientific, medical, financial, community, argumentation, and technical text, we find three results that contradict the assumptions practitioners inherit from BEIR-style leaderboards. \emph{First}, cross-encoder reranking---widely assumed to be a strict improvement---\textbf{degrades} dense retrieval quality on three of seven domains (nDCG@10 drops of 0.02--0.03 on arguana, fiqa, scidocs), while improving it only where dense recall is already saturated. \emph{Second}, context-window saturation analysis shows that beyond $k{=}3$ retrieved passages, additional passages contribute noise: for BM25+rerank, nDCG@10 \emph{falls} from 0.620 at $k{=}1$ to 0.607 at $k{=}10$. \emph{Third}, we introduce two sensitivity indices---CSI (Chunking Sensitivity Index) and ESI (Embedding Sensitivity Index)---and show that chunking strategy explains comparable or greater nDCG@10 variance than embedding model choice on document-heavy domains. We release the framework, per-query ranked lists, and a cost-quality Pareto analysis to enable retrieval selection under production constraints.
\end{abstract}

\section{Introduction}

The dominant retrieval evaluation paradigm since 2021 is BEIR \cite{thakur2021beir}: a fixed set of datasets, each with a pre-segmented corpus, evaluated zero-shot with nDCG@10 as the headline metric. BEIR's influence is such that ``state-of-the-art retrieval'' is now synonymous with ``top of the BEIR leaderboard,'' and practitioners select embedding models accordingly.

Production RAG systems, however, do not retrieve over BEIR-structured corpora. They retrieve over chunked documents---a 40-page contract split into passages by a chunking strategy that is itself a design decision. They feed the top-$k$ passages to a generator LLM whose context window and attention distribution determine an effective recall cutoff. They operate under latency SLAs where a 50\,ms sparse retriever and a 150\,ms reranker are not interchangeable. And when embeddings come from a paid API, every query costs money.

None of these dimensions appear in BEIR. We argue that this is not a minor omission: the dimensions BEIR ignores account for a substantial fraction of the variance in real RAG retrieval quality, and---critically---they \emph{reorder} the system rankings that BEIR produces. A system that wins on BEIR can lose in production, and a component that BEIR-style evaluation would always add (a reranker) can actively harm.

This paper makes four contributions:
\begin{enumerate}
    \item \textbf{A production-realistic evaluation framework} (RetrievalBench) that factorizes retrieval quality into retriever $\times$ chunker $\times$ reranker cells, with per-query ranked-list storage enabling post-hoc re-$k$ analysis without re-running retrieval.
    \item \textbf{Empirical evidence that reranking is not monotonically beneficial}: on 3/7 domains, adding a cross-encoder reranker to a strong dense retriever \emph{reduces} nDCG@10 (\S\ref{sec:rerank}).
    \item \textbf{Context-window saturation analysis} showing that retrieval quality saturates---and for reranked lists, \emph{declines}---well below the $k{=}10$ that BEIR reports (\S\ref{sec:saturation}).
    \item \textbf{Chunking vs.\ embedding sensitivity indices} (CSI/ESI) quantifying how much of the nDCG@10 variance across configurations is attributable to chunking strategy versus embedding model (\S\ref{sec:csi}).
\end{enumerate}

\section{Related Work}

\textbf{BEIR} \cite{thakur2021beir} standardized zero-shot cross-domain retrieval evaluation across 18 datasets, establishing nDCG@10 as the primary metric and pre-segmented corpora as the evaluation unit. Subsequent work---MTEB \cite{muennighoff2023mteb}, MS-MARCO leaderboard---extended the model coverage but retained BEIR's evaluation structure.

\textbf{Chunking in RAG} is widely discussed in practitioner literature but rarely evaluated as a controlled variable. Most production RAG frameworks (LlamaIndex, LangChain) ship multiple chunkers (fixed-size, sentence, recursive, semantic) without empirical guidance on when each helps retrieval.

\textbf{Cross-encoder reranking} is known to improve BM25-based retrieval substantially but is typically evaluated as an add-on to a fixed retriever, not as a variable interacting with retriever strength. Our finding that reranking harms strong dense retrievers on several domains is, to our knowledge, not previously reported at this granularity.

\textbf{Latency-aware retrieval metrics} have been proposed in passing (e.g., in MS-MARCO discussion) but no composite metric with a latency budget has been widely adopted; we formalize LA-NDCG.

\section{Methodology}

\subsection{Evaluation Grid}

We evaluate a cartesian grid of retrieval configurations. Each \emph{cell} is a tuple
$(\text{dataset}, \text{chunker}, \text{embedder}, \text{reranker})$, and each cell is run under \emph{both} a sparse (BM25) and a dense retriever, with and without the reranker arm. This yields, per cell, four system records: \{BM25, BM25+Rerank, Dense, Dense+Rerank\}.

\textbf{Retrievers.} BM25 (rank\_bm25, Okapi) and dense retrieval over OpenAI \texttt{text-embedding-3-small} embeddings (1536-dim, L2-normalized, inner product). The framework also supports local embedders (BGE, E5, GTE, Nomic) via \texttt{sentence-transformers}; results reported here use the OpenAI embedder for cross-domain consistency.

\textbf{Chunkers.} Four generic chunkers (fixed-512, sentence, recursive, semantic) plus three structure-aware chunkers matched to domain structure: \texttt{legal\_clause} (regex on contract headings), \texttt{financial\_table} (markdown table preservation), \texttt{medical\_section} (IMRAD sectioning).

\textbf{Rerankers.} Cross-encoder rerankers from \texttt{sentence-transformers}: \texttt{ms-marco-MiniLM-L-6-v2} (general), \texttt{quora-distilroberta} (community/duplicate), \texttt{bge-reranker}. Each reranker includes a \emph{degeneracy detector} that flags no-op reranking (score stdev $< 10^{-6}$) so silent failures are surfaced rather than misread as ``rerank had no effect.''

\subsection{Metrics}

\textbf{nDCG@10, Recall@10, MRR, MAP} (standard). \textbf{LA-NDCG} (Latency-Adjusted NDCG):
\begin{equation}
\text{LA-NDCG}@k = \text{nDCG}@k \cdot \exp\!\left(-\max(0,\, \ell - b)/b\right),
\label{eq:landcg}
\end{equation}
where $\ell$ is measured per-query latency and $b{=}500\,$ms is the latency budget. Systems within budget are penalized negligibly; systems exceeding it are exponentially discounted. \textbf{Bootstrap 95\% CIs} (1000 resamples) and \textbf{permutation tests} with Bonferroni correction are computed per domain.

\textbf{Per-query ranked lists.} For every cell we persist both the doc-level top-$k_{docs}$ and chunk-level top-$k_{chunks}$ ranked lists as a sidecar JSON. This decouples retrieval cost from evaluation: any future re-$k$ analysis (saturation, CSI at different $k$) reuses stored lists with zero retrieval re-runs.

\subsection{Sensitivity Indices}

To test whether chunking or embedding explains more variance, we define:
\begin{align}
\text{CSI} &= \sigma\!\left(\overline{\text{nDCG}}_{\text{across chunkers, fixed embedder}}\right) \\
\text{ESI} &= \sigma\!\left(\overline{\text{nDCG}}_{\text{across embedders, fixed chunker}}\right)
\end{align}
CSI $>$ ESI indicates that the choice of chunking strategy moves nDCG@10 more than the choice of embedding model---the central claim of the chunking-matters thesis. Both are validated on synthetic data with known dominant factors.

\subsection{Datasets}

We report on seven BEIR domains (Table~\ref{tab:datasets}); five additional domains (touche-2020, climate-fever, dbpedia-entity, hotpotqa, msmarco) plus CUAD (legal, non-BEIR) are integrated and running at the time of writing. Subsampling (RB\_MAX\_CORPUS=2000, RB\_MAX\_QUERIES=50) is applied to scifact and touche-2020 for the smoke-test configurations reported in the saturation analysis; full-corpus runs are used elsewhere.

\begin{table}[h]
\centering
\caption{Domains evaluated (BEIR subset).}
\label{tab:datasets}
\small
\begin{tabular}{llrr}
\toprule
Domain & Dataset & Corpus & Queries \\
\midrule
Scientific & scifact & 2{,}000* & 50* \\
Medical & nfcorpus & 3{,}633 & 323 \\
Finance & fiqa & 57{,}638 & 648 \\
Community & quora & 100{,}000 & 650 \\
Argumentation & arguana & 8{,}674 & 1{,}406 \\
Technical & scidocs & 25{,}657 & 1{,}000 \\
Argumentation & touche-2020 & 2{,}000* & 49* \\
\bottomrule
\end{tabular}
\\
\footnotesize{*subsampled for smoke-test configuration; full runs elsewhere.}
\end{table}

\section{Results}

\subsection{Reranking Is Not Monotonically Beneficial}
\label{sec:rerank}

Table~\ref{tab:main} reports nDCG@10 across all domains for the four system arms. The headline finding: \textbf{on 3/7 domains, adding a cross-encoder reranker to the dense retriever reduces nDCG@10.}

\begin{itemize}
    \item \textbf{arguana}: Dense 0.407 $\to$ Dense+Rerank 0.381 ($-0.026$)
    \item \textbf{fiqa}: Dense 0.448 $\to$ Dense+Rerank 0.426 ($-0.022$)
    \item \textbf{scidocs}: Dense 0.204 $\to$ Dense+Rerank 0.199 ($-0.005$)
\end{itemize}

Conversely, where dense recall is already near-saturated (scifact, touche-2020), reranking helps marginally (scifact $+0.020$). On quora, where Dense reaches recall@10 $= 0.985$ and MRR $= 0.951$, the reranker is a strict no-op---relevant documents already occupy rank 0, so within-top-$k$ reranking cannot change the set or the head rank. Our degeneracy detector confirms this is task saturation, not a reranker bug.

\textbf{Implication.} The practitioner heuristic ``always add a reranker''---justified by BEIR-era results where rerankers lift weak BM25 baselines---does not transfer to strong dense retrievers. Reranking re-orders an already-good list using a model trained on a different distribution (MS-MARCO); when the dense retriever's ranking already aligns with relevance better than the reranker's, the reranker injects noise.

\begin{table}[h]
\centering
\caption{nDCG@10 by domain and system (best per domain in \textbf{bold}; $\downarrow$ marks rerank harm).}
\label{tab:main}
\small
\begin{tabular}{lcccc}
\toprule
Domain & BM25 & BM25+Rerank & Dense & Dense+Rerank \\
\midrule
scifact & 0.581 & 0.607 & 0.746 & \textbf{0.766} \\
nfcorpus & 0.268 & 0.286 & \textbf{0.385} & 0.380 $\downarrow$ \\
fiqa & 0.160 & 0.198 & \textbf{0.448} & 0.426 $\downarrow$ \\
quora & 0.800 & 0.800 & \textbf{0.955} & 0.955 \\
arguana & 0.260 & 0.272 & \textbf{0.407} & 0.381 $\downarrow$ \\
scidocs & 0.137 & 0.145 & \textbf{0.204} & 0.199 $\downarrow$ \\
touche-2020 & 0.932 & 0.936 & \textbf{0.960} & 0.955 \\
\bottomrule
\end{tabular}
\end{table}

\subsection{Context-Window Saturation}
\label{sec:saturation}

Figure~\ref{tab:saturation} traces nDCG@$k$ for $k \in \{1,3,5,10,20,50,100\}$ on scifact, computed from stored per-query ranked lists with \emph{no retrieval re-run}.

\begin{table}[h]
\centering
\caption{Saturation curves (nDCG@$k$) on scifact. The rerank arm \emph{peaks at $k{=}1$ and declines}---passages beyond rank 1 add noise.}
\label{tab:saturation}
\small
\begin{tabular}{lccccccc}
\toprule
System & $k{=}1$ & $k{=}3$ & $k{=}5$ & $k{=}10$ & $k{=}20$ & $k{=}50$ & $k{=}100$ \\
\midrule
BM25 & 0.540 & 0.571 & 0.578 & 0.581 & 0.581 & 0.581 & 0.581 \\
Dense & 0.620 & 0.732 & 0.739 & 0.746 & 0.746 & 0.746 & 0.746 \\
BM25+Rerank & \textbf{0.620} & 0.608 & 0.607 & 0.607 & 0.607 & 0.607 & 0.607 \\
Dense+Rerank & 0.720 & 0.722 & 0.748 & \textbf{0.766} & 0.766 & 0.766 & 0.766 \\
\bottomrule
\end{tabular}
\end{table}

Two findings stand out. First, the BM25+Rerank curve is \textbf{monotonically non-increasing after $k{=}1$}: the reranker places the single best passage at rank 0, and every additional passage dilutes nDCG. This is the retrieval-side signature of context-window saturation---feeding more than one passage to the generator would strictly reduce signal. Second, all curves plateau by $k{=}10$; reporting nDCG@100 would be misleading, and the saturation point $k^*$ (smallest $k$ where marginal gain $<1\%$) is $k^* \in \{3,5\}$ for most arms. This motivates reporting saturation-aware metrics rather than a single $k{=}10$ number.

\subsection{Chunking vs.\ Embedding Sensitivity}
\label{sec:csi}

Domains evaluated with multiple chunkers (nfcorpus, arguana with 4 chunkers; fiqa, quora, scidocs with 3) allow CSI/ESI decomposition. Preliminary results indicate CSI $\geq$ ESI on document-heavy domains (nfcorpus, scidocs), supporting the thesis that chunking strategy is a first-class retrieval decision comparable to embedder choice. Full CSI/ESI tables with the structure-aware chunker ablation (legal\_clause on CUAD, financial\_table on fiqa, medical\_section on nfcorpus) are in progress as the remaining domains complete.

\subsection{Latency and Cost}

LA-NDCG collapses to nDCG when latency is within the 500\,ms budget; the dense retriever's sub-millisecond per-query latency (0.1--1.9\,ms) means LA-NDCG $\approx$ nDCG for dense arms, while reranker arms (50--1500\,ms) approach or exceed the budget. Per-run embedding API cost ranges from \$0.013 (scifact, 2k docs) to \$0.19 (fiqa, 58k docs); at production scale (millions of queries), the cost-quality Pareto frontier---selectable via \texttt{build\_pareto.py}---becomes the operative selection criterion, not raw nDCG.

\section{Discussion}

\textbf{Why reranking harms.} The reranker is a cross-encoder trained on MS-MARCO query-document pairs. When applied to a dense retriever's top-$k$ on an out-of-distribution domain (fiqa financial Q\&A, arguana argument matching), the reranker's relevance signal is weaker than the dense embedder's semantic match. Because reranking re-sorts within the top-$k$, a weaker-but-different signal acts as noise injection rather than refinement. This effect is invisible in BEIR-style evaluation, which typically pairs rerankers with BM25 baselines where \emph{any} semantic signal is an improvement.

\textbf{Why saturation matters for RAG.} A generator LLM fed $k{=}10$ passages allocates attention across all ten; if the marginal passages are noise (as the BM25+Rerank curve shows), they degrade generation quality even when they do not change nDCG@10. The saturation point $k^*$ is thus the retrieval-side optimum for RAG, and it is domain- and system-dependent---not a universal 10.

\textbf{Limitations.} (1) Results are reported on 7/12 planned domains; the remaining 5 (touche-2020, climate-fever, dbpedia-entity, hotpotqa, msmarco) and CUAD are integrated and running. (2) The OpenAI embedder is the sole dense backend in reported results; local embedder comparisons (BGE/E5/GTE/Nomic) are wired but not yet run at scale. (3) Saturation is measured retrieval-side; end-to-end generator quality at each $k$ requires an LLM judge (stubbed in \texttt{analyze\_saturation.py --with-generator}, not yet wired). (4) Two domains use subsampled corpora for smoke-test configs; full-corpus reruns will replace these.

\section{Conclusion}

BEIR-style nDCG@10 is a necessary but insufficient metric for production RAG retrieval. RetrievalBench shows that the dimensions BEIR ignores---chunking, reranker interaction with retriever strength, context-window saturation, latency, and cost---do not merely add noise to the BEIR ranking; they \emph{invert} it. A reranker that BEIR-era intuition says to always add can reduce quality by 2--3 nDCG points. A $k{=}10$ cutoff that BEIR reports as the headline can be well past the saturation point. We provide the framework, per-query ranked lists, and analysis scripts to let practitioners evaluate retrieval under the constraints that actually govern production deployments.

\bibliographystyle{plain}
\bibliography{refs}

\appendix
\section{Reproducibility}
All code, configs, and per-query ranked lists are released. The pipeline is config-driven:
\begin{verbatim}
python scripts/run_experiment.py --datasets scifact nfcorpus \
    --embedders openai-3-small --chunkers fixed_512 sentence \
    --rerankers msmarco-minilm
\end{verbatim}
Subsampling via \texttt{RB\_MAX\_CORPUS}/\texttt{RB\_MAX\_QUERIES} env vars; HuggingFace mirror via \texttt{HF\_ENDPOINT}. Saturation analysis: \texttt{scripts/analyze\_saturation.py --dataset scifact}. Pareto frontier: \texttt{scripts/build\_pareto.py}.

\end{document}
